% ================================================================
% V3 Stage: Causal Utilization and Spatiotemporal Dynamics Complete Summary Tables
% ================================================================

\begin{table}[htbp]
\centering
\small
\caption{V3 State-Induction Gate 判定結果：因果介入パイプライン継続可否の評価規約。各判定軸の信頼区間下限／上限が 1--9 raw scale 閾値（Sufficiency: 0.10, Specificity: 0.05, Endogenous: 0.05）および TVD 上限（0.15）を満たすかを判定。}
\label{tab:v3_gate_decision}
\begin{tabular}{l cccc cccc c}
\toprule
\textbf{Axis} & \multicolumn{2}{c}{\textbf{Sufficiency ($\beta_1$)}} & \multicolumn{2}{c}{\textbf{Specificity ($\beta_1$)}} & \multicolumn{2}{c}{\textbf{Endogenous ($M$)}} & \multicolumn{2}{c}{\textbf{Topic Control (TVD)}} & \textbf{Gate Decision} \\
\cmidrule(lr){2-3} \cmidrule(lr){4-5} \cmidrule(lr){6-7} \cmidrule(lr){8-9}
 & \textbf{CI Low} & \textbf{Pass?} & \textbf{CI Low} & \textbf{Pass?} & \textbf{CI Low} & \textbf{Pass?} & \textbf{CI High} & \textbf{Pass?} & \\
\midrule
Valence    & 4.34e-04   & $\times$ & $-$0.0025  & $\times$ & $-$0.0366  & $\times$ & 0.0043     & \checkmark & \textbf{NO\_GO} \\
Arousal    & -5.13e-05  & $\times$ & -4.18e-04  & $\times$ & 0.0013     & $\times$ & 0.0045     & \checkmark & \textbf{NO\_GO} \\
\bottomrule
\end{tabular}
\vspace{1ex}
\begin{minipage}{\linewidth}
\footnotesize
\textbf{Note:} State-Induction Gate 判定基準（Sufficiency: $\text{CI}_{\text{low}} > 0.10$, Specificity: $\text{CI}_{\text{low}} > 0.05$, Endogenous Relevance: $\text{CI}_{\text{low}} > 0.05$, Topic TVD: $\text{CI}_{\text{high}} < 0.15$）。全軸でSufficiency / Specificity / Endogenousの事前定義thresholdを満たさなかったため、パイプライン判定は \textbf{NO\_GO}（\texttt{pipeline\_continues = False}）となり、以降の分析は探索的・定性的分析として位置づけられる。
\end{minipage}
\end{table}

\begin{table}[htbp]
\centering
\small
\caption{V3 探索的時空間ダイナミクス（Qwen 2.5 1.5B Discovery）：直接読み出しピーク深度 $d_D$ と因果介入ピーク深度 $d_C$ の空間的解離、および下流層媒介減衰率。}
\label{tab:v3_spatiotemporal_dynamics}
\begin{tabular}{l ccccc ccc}
\toprule
 & \multicolumn{5}{c}{\textbf{Peak \& Center-of-Mass Depth}} & \multicolumn{3}{c}{\textbf{Mediated Attenuation}} \\
\cmidrule(lr){2-6} \cmidrule(lr){7-9}
\textbf{Axis} & $d_D^{\text{peak}}$ & $d_C^{\text{peak}}$ & $\Delta d^{\text{peak}}$ & $d_D^{\text{center}}$ & $d_C^{\text{center}}$ & Layer ($d$) & Atten. Ratio & 95\% CI \\
\midrule
Valence    & 0.852    & 0.741    & $-$0.111 & 0.737    & 0.609    & L3 (0.11)    & 0.87\%       & [-1.77, 3.42]\% \\
Arousal    & 0.852    & 0.778    & $-$0.074 & 0.717    & 0.756    & L3 (0.11)    & 0.58\%       & [0.01, 1.13]\% \\
\bottomrule
\end{tabular}
\vspace{1ex}
\begin{minipage}{\linewidth}
\footnotesize
\textbf{Note:} 全軸で $\Delta d^{\text{peak}} < 0$ であり、直接読み出しが可能な表現層よりも前の層で因果効果が最大化する「時空間的解離」が観測された。しかし、媒介減衰率はいずれも微小（$< 1.0\%$）であり、下流媒介仮説の成立は限定的である。
\end{minipage}
\end{table}

\begin{table}[htbp]
\centering
\small
\caption{V3 探索的媒介減衰詳細（Mediated Attenuation Analysis）：中間層（Mediator Layer）への介入による総変位量 $T$、残差変位 $R$、絶対減衰量 $M = T - R$、ランダム統制 $M_{\text{rand}}$、正味媒介減衰量 $M_{\text{net}} = M - M_{\text{rand}}$、減衰率（Attenuation Ratio）、およびブートストラップ95\%信頼区間。}
\label{tab:v3_mediated_attenuation}
\begin{tabular}{ll cccccc cc}
\toprule
\textbf{Axis} & \textbf{Mediator ($d$)} & $T$ & $R$ & $M$ & $M_{\text{rand}}$ & $M_{\text{net}}$ & \textbf{Atten. Ratio} & \textbf{95\% CI ($M$)} & \textbf{95\% CI ($M_{\text{net}}$)} \\
\midrule
Valence    & L3 (0.11)        & ---      & ---      & 9.01e-04   & ---        & ---        & 0.87\%         & [$-$0.0011, 0.0029]    & --- \\
Arousal    & L3 (0.11)        & ---      & ---      & 2.71e-04   & ---        & ---        & 0.58\%         & [-1.53e-04, 6.57e-04]  & --- \\
\bottomrule
\end{tabular}
\vspace{1ex}
\begin{minipage}{\linewidth}
\footnotesize
\textbf{Note:} $M = T - R$ は媒介層の潜在表現統制による変位減衰量、$M_{\text{net}} = M - M_{\text{rand}}$ はランダム部分空間統制を差し引いた正味媒介減衰量。ValenceおよびArousalともに 95\% CI がゼロを跨ぎ、統計的に堅牢な媒介効果は支持されない。
\end{minipage}
\end{table}

\begin{table}[htbp]
\centering
\small
\caption{V3 cross-family replication details。Llama 3.2、Gemma 3、OLMo 2におけるH1--H4のValence / Arousal別推定値、family-specific 95\% confidence interval、replication criterion、およびPass / Failを示す。H1のdirectionはQwen Discovery終了後、replication-family outcomesを評価する前にfreezeした。}
\label{tab:v3_confirmatory_details}
\begin{tabular}{ll cccc c}
\toprule
\textbf{Family} & \textbf{Hypothesis} & \textbf{Axis} & \textbf{Estimate} & \textbf{Family 95\% CI} & \textbf{Threshold} & \textbf{Pass/Fail} \\
\midrule
\multirow{12}{*}{\textbf{Llama 3.2}} & H1 Peak Dissoc. ($\Delta d^*$)         & Valence  & $-$0.87    & [$-$1.00, $-$0.27] & $\text{CI}_{\text{high}} < 0$ & \checkmark \textbf{PASS} \\
                             & H1 Peak Dissoc. ($\Delta d^*$)         & Arousal  & 0.07       & [$-$0.47, 0.40]    & $\text{CI}_{\text{high}} < 0$ & $\times$ FAIL \\
                             & H1 Center Dissoc. ($\Delta\bar{d}$)    & Valence  & $-$0.32    & [$-$0.52, $-$0.11] & $\text{CI}_{\text{high}} < 0$ & \checkmark \textbf{PASS} \\
                             & H1 Center Dissoc. ($\Delta\bar{d}$)    & Arousal  & $-$0.12    & [$-$0.53, $-$0.09] & $\text{CI}_{\text{low}} > 0$ & $\times$ FAIL \\
                             & H2: Sufficiency ($\beta_1$)            & Valence  & 0.0032     & [0.0022, 0.0042]   & $\text{CI}_{\text{low}} > 0.10$ & $\times$ FAIL \\
                             & H2: Sufficiency ($\beta_1$)            & Arousal  & 0.0018     & [0.0015, 0.0022]   & $\text{CI}_{\text{low}} > 0.10$ & $\times$ FAIL \\
                             & H3: Mediation ($M$)                    & Valence  & 6.87e-05   & [$-$0.0028, 0.0027] & $\text{CI}_{\text{low}}(M) > 0$ & $\times$ FAIL \\
                             & H3: Mediation ($M$)                    & Arousal  & 0.0013     & [2.25e-04, 0.0025] & $\text{CI}_{\text{low}}(M) > 0$ & \checkmark \textbf{PASS} \\
                             & H3: Net vs. Random ($M_{\text{net}}$)  & Valence  & 3.15e-04   & [$-$0.0019, 0.0024] & $\text{CI}_{\text{low}}(M_{\text{net}}) > 0$ & $\times$ FAIL \\
                             & H3: Net vs. Random ($M_{\text{net}}$)  & Arousal  & 9.63e-04   & [-1.85e-04, 0.0022] & $\text{CI}_{\text{low}}(M_{\text{net}}) > 0$ & $\times$ FAIL \\
                             & H4: Contrast ($\Delta C$)              & Valence  & 0.0445     & [0.0415, 0.0473]   & $\text{CI}_{\text{low}} > 0$ & \checkmark \textbf{PASS} \\
                             & H4: Contrast ($\Delta C$)              & Arousal  & 0.0656     & [0.0627, 0.0684]   & $\text{CI}_{\text{low}} > 0$ & \checkmark \textbf{PASS} \\
\midrule
\multirow{12}{*}{\textbf{Gemma 3}} & H1 Peak Dissoc. ($\Delta d^*$)         & Valence  & 0.00       & [$-$0.64, 0.36]    & $\text{CI}_{\text{high}} < 0$ & $\times$ FAIL \\
                             & H1 Peak Dissoc. ($\Delta d^*$)         & Arousal  & $-$0.28    & [$-$1.00, 0.00]    & $\text{CI}_{\text{high}} < 0$ & $\times$ FAIL \\
                             & H1 Center Dissoc. ($\Delta\bar{d}$)    & Valence  & $-$0.13    & [$-$0.23, 0.38]    & $\text{CI}_{\text{high}} < 0$ & $\times$ FAIL \\
                             & H1 Center Dissoc. ($\Delta\bar{d}$)    & Arousal  & $-$0.08    & [$-$0.34, 0.21]    & $\text{CI}_{\text{low}} > 0$ & $\times$ FAIL \\
                             & H2: Sufficiency ($\beta_1$)            & Valence  & $-$0.0011  & [$-$0.0036, 0.0020] & $\text{CI}_{\text{low}} > 0.10$ & $\times$ FAIL \\
                             & H2: Sufficiency ($\beta_1$)            & Arousal  & -9.01e-04  & [$-$0.0020, 6.86e-05] & $\text{CI}_{\text{low}} > 0.10$ & $\times$ FAIL \\
                             & H3: Mediation ($M$)                    & Valence  & 0.0031     & [$-$0.0088, 0.0141] & $\text{CI}_{\text{low}}(M) > 0$ & $\times$ FAIL \\
                             & H3: Mediation ($M$)                    & Arousal  & -2.29e-05  & [$-$0.0041, 0.0042] & $\text{CI}_{\text{low}}(M) > 0$ & $\times$ FAIL \\
                             & H3: Net vs. Random ($M_{\text{net}}$)  & Valence  & 0.0086     & [-7.03e-04, 0.0179] & $\text{CI}_{\text{low}}(M_{\text{net}}) > 0$ & $\times$ FAIL \\
                             & H3: Net vs. Random ($M_{\text{net}}$)  & Arousal  & 2.64e-04   & [$-$0.0033, 0.0036] & $\text{CI}_{\text{low}}(M_{\text{net}}) > 0$ & $\times$ FAIL \\
                             & H4: Contrast ($\Delta C$)              & Valence  & 0.0059     & [0.0020, 0.0095]   & $\text{CI}_{\text{low}} > 0$ & \checkmark \textbf{PASS} \\
                             & H4: Contrast ($\Delta C$)              & Arousal  & 0.0082     & [0.0064, 0.0098]   & $\text{CI}_{\text{low}} > 0$ & \checkmark \textbf{PASS} \\
\midrule
\multirow{12}{*}{\textbf{OLMo 2}} & H1 Peak Dissoc. ($\Delta d^*$)         & Valence  & $-$0.67    & [$-$1.00, $-$0.07] & $\text{CI}_{\text{high}} < 0$ & \checkmark \textbf{PASS} \\
                             & H1 Peak Dissoc. ($\Delta d^*$)         & Arousal  & $-$0.93    & [$-$0.93, 0.27]    & $\text{CI}_{\text{high}} < 0$ & $\times$ FAIL \\
                             & H1 Center Dissoc. ($\Delta\bar{d}$)    & Valence  & $-$0.24    & [$-$0.44, $-$0.04] & $\text{CI}_{\text{high}} < 0$ & \checkmark \textbf{PASS} \\
                             & H1 Center Dissoc. ($\Delta\bar{d}$)    & Arousal  & $-$0.62    & [$-$0.63, $-$0.10] & $\text{CI}_{\text{low}} > 0$ & $\times$ FAIL \\
                             & H2: Sufficiency ($\beta_1$)            & Valence  & 4.27e-05   & [-8.42e-04, 9.50e-04] & $\text{CI}_{\text{low}} > 0.10$ & $\times$ FAIL \\
                             & H2: Sufficiency ($\beta_1$)            & Arousal  & -1.43e-04  & [-3.28e-04, 6.10e-05] & $\text{CI}_{\text{low}} > 0.10$ & $\times$ FAIL \\
                             & H3: Mediation ($M$)                    & Valence  & 1.28e-04   & [$-$0.0024, 0.0028] & $\text{CI}_{\text{low}}(M) > 0$ & $\times$ FAIL \\
                             & H3: Mediation ($M$)                    & Arousal  & -1.37e-05  & [-7.29e-04, 6.06e-04] & $\text{CI}_{\text{low}}(M) > 0$ & $\times$ FAIL \\
                             & H3: Net vs. Random ($M_{\text{net}}$)  & Valence  & -2.66e-04  & [$-$0.0021, 0.0015] & $\text{CI}_{\text{low}}(M_{\text{net}}) > 0$ & $\times$ FAIL \\
                             & H3: Net vs. Random ($M_{\text{net}}$)  & Arousal  & -1.86e-05  & [-4.88e-04, 4.11e-04] & $\text{CI}_{\text{low}}(M_{\text{net}}) > 0$ & $\times$ FAIL \\
                             & H4: Contrast ($\Delta C$)              & Valence  & 0.0179     & [0.0161, 0.0199]   & $\text{CI}_{\text{low}} > 0$ & \checkmark \textbf{PASS} \\
                             & H4: Contrast ($\Delta C$)              & Arousal  & 0.0678     & [0.0654, 0.0704]   & $\text{CI}_{\text{low}} > 0$ & \checkmark \textbf{PASS} \\
\bottomrule
\end{tabular}
\vspace{1ex}
\begin{minipage}{\linewidth}
\footnotesize
\textbf{Note:} 各ファミリー固有の95\%ブートストラップ信頼区間および判定基準（H1: Qwen Discoveryと同方向の空間的解離 $\text{CI}_{\text{high}} < 0$ または $\text{CI}_{\text{low}} > 0$; H2: $\text{CI}_{\text{low}} > 0.10$; H3: $\text{CI}_{\text{low}}(M) > 0$ かつ $\text{CI}_{\text{low}}(M_{\text{net}}) > 0$; H4: $\text{CI}_{\text{low}} > 0$）に基づく評価。H1では、Qwen Discoveryで観測されたspatial dissociationのdirectionを、Llama、Gemma、OLMoのoutcomeを評価する前にfreezeし、同方向のeffectについてfamily-specific bootstrap CIが0を除外するかを評価した。H1のdirectional dissociationはValenceにおいてLlama 3.2およびOLMo 2で再現された（PASS）が、Arousalでは再現されず、3 families全体での一貫した解離は支持されなかった。H3の支持には内因性変位減衰量 $M$ の信頼区間下限が正であること（$\text{CI}_{\text{low}}(M) > 0$）に加え、ランダム部分空間統制を差し引いた正味減衰量 $M_{\text{net}}$ の信頼区間下限も正であること（$\text{CI}_{\text{low}}(M_{\text{net}}) > 0$）が要求される。H2およびH3はいずれのモデル・軸でもreplication基準を満たさなかった（FAIL）。H4では各モデルで軸レベルの正の効果量（$\text{CI}_{\text{low}} > 0$）が観測されたが、Gate判定がNO\_GOであるためパイプライン全体のConfirmationは不成立となり、探索的証拠として位置づけられる。
\end{minipage}
\end{table}

\begin{table}[htbp]
\centering
\small
\caption{V3 独立ファミリー再現性マトリックス：探索段階（Qwen）で得られた知見の独立3ファミリー（Llama 3.2, Gemma 3, OLMo 2）における検証結果総括。各仮説はValenceおよびArousalの双方が基準を満たした場合のみ支持（$\checkmark$）とされる。}
\label{tab:v3_confirmatory_matrix}
\begin{tabular}{l cccc c}
\toprule
\textbf{Family} & \textbf{H1 (Dissociation)} & \textbf{H2 (Sufficiency)} & \textbf{H3 (Mediation)} & \textbf{H4 (Temporal Contrast)} & \textbf{All H1--H4 Criteria?} \\
\midrule
Llama 3.2      & $\times$     & $\times$     & $\times$     & \checkmark   & \textbf{NO} \\
Gemma 3        & $\times$     & $\times$     & $\times$     & \checkmark   & \textbf{NO} \\
OLMo 2         & $\times$     & $\times$     & $\times$     & \checkmark   & \textbf{NO} \\
\bottomrule
\end{tabular}
\vspace{1ex}
\begin{minipage}{\linewidth}
\footnotesize
\textbf{Note:} 全4仮説の判定マトリックス。各セルは該当仮説についてValenceおよびArousalの双方が事前定義された95\% CI criterionを満たした場合に$\checkmark$ とする。H1のfamily-level criterionはValence / Arousal双方のaxis-level criterionを要求するため、いずれのfamilyでも成立しなかった。ただしValenceのdirectional criterionはLlama 3.2およびOLMo 2で満たされた。H2およびH3は、いずれのfamilyでもfamily-level criterionを満たさなかった。H4のtemporal contrastは3 familyすべてでValence / Arousal双方がaxis-level criterionを満たした。ただし、事前のState Induction GateがNO\_GOであったため、V3全体としてfamily-general causal mechanismが確立されたとは解釈しない（All Criteria = NO）。
\end{minipage}
\end{table}
